\RequirePackage{pdfmanagement}
\SetKeys[document/metadata]{pdfversion=1.7, pdfstandard=A-2u, lang=en-GB}
\input{preamble/pdfA-compliance/glyphtounicode}
\documentclass[11pt,a4paper]{report}

\newcommand*{\PaperTitle}{Title, math allowed: \(W^a(3)\ge1\)}
\newcommand*{\PaperTitlePlain}{Title in plain text for the PDF metadata}
\newcommand*{\PaperShortTitle}{Running head}
\newcommand*{\PaperSubtitle}{chemistry-open-problems \textperiodcentered{} DD Month YYYY \textperiodcentered{} preliminary}
\newcommand*{\PaperCredits}{Solver: Model (effort ...) \textperiodcentered{} Independent verification: Model (effort ...) or pending}
\newcommand*{\PaperKeywords}{keyword, keyword}
\input{chemop-preamble}
\addbibresource{references.bib}

\begin{document}
\makepapertitle

\begin{paperabstract}
One paragraph: the open question, who asked it, the answer, the nature of the proof or certificate.
\end{paperabstract}

\section*{The question}
Source and verbatim open statement:
\begin{openquote}
``Quoted sentence that makes the problem open.''
\end{openquote}

\begin{theorem}\label{thm:main}
Statement.
\end{theorem}

\section{Proof}
\begin{lemma}[note=optional name]\label{lem:one}
Statement.
\end{lemma}
\begin{proof}
Proof. Cite with \autocite{Key}; refer with \zcref{thm:main}.
\end{proof}

\section{Computer verification}
What each script in \texttt{scripts/} checks, and how (exact arithmetic, intervals, solver certificates).

\paperreferences
\end{document}
